% cap_resultados.tex
\chapter{Resultados e Discussão}
\label{chap:resultados}

Este capítulo apresenta os resultados do procedimento descrito no Capítulo~\ref{chap:materiais_metodos}. Todas as tabelas e figuras numéricas são geradas por programa a partir das previsões salvas de cada execução e incluídas sem transcrição manual. Os resultados são apresentados em dois períodos. O \emph{período de confirmação} (2021 a 2023) não foi consultado em nenhuma decisão de projeto e é a evidência principal, por isso as tabelas do texto se referem a ele. O \emph{teste histórico} (27 de maio a 31 de dezembro de 2020) já foi consultado durante o desenvolvimento, e suas tabelas estão no Apêndice~\ref{ap:teste_historico}; sempre que os dois períodos discordam, isso é indicado.

Nas tabelas de comparação, a diferença de MAE é calculada como o MAE do modelo menos o MAE da referência; valores negativos indicam que o modelo erra menos. O intervalo de confiança (IC) de 95\% vem do bootstrap em blocos e o valor-\(p\) (DM) vem do teste de Diebold-Mariano (Seção~\ref{sec:inferencia_estatistica}). A coluna ``Veredito'' resume a comparação: ``melhor'' significa que o IC está inteiramente abaixo de zero (o modelo erra menos), ``pior'' que está inteiramente acima, e ``empate estatístico'' que o IC contém o zero, isto é, não há evidência de diferença. As médias de MAE da comparação principal são de quatro sementes.

\section{Comparação Principal}
\label{sec:comparacao_principal}

A Tabela~\ref{tab:metricas_holdout} apresenta as métricas de erro de cada modelo e das referências no período de confirmação, com a mesma preparação, o mesmo orçamento de busca e a mesma perda de Huber (o XGBoost usa o erro absoluto). A última coluna é a habilidade em relação à média móvel de 24 horas: um valor de 0,10 significa erro 10\% menor que o dessa referência.

\inputtabela{metricas_holdout}

O XGBoost obteve o menor erro em todas as medidas globais: MAE de 2,767~\(\mu\)g/m\(^3\), RMSE de 4,072 e \(R^2\) de 0,640. Em seguida vêm a regressão Ridge (MAE de 2,794), as duas versões do Seq2Seq com atenção (2,870 e 2,876), a LSTM recursiva (2,944), o Seq2Seq básico (2,965) e a LSTM direta (2,979). Entre os modelos treinados, a diferença entre o melhor e o pior é de 0,212~\(\mu\)g/m\(^3\), cerca de 7,7\% do MAE do melhor. Todos superam as referências simples: a média móvel de 24 horas tem MAE de 3,072, e a persistência (repetir o último valor), de 3,831. O XGBoost, por exemplo, tem erro 9,9\% menor que a média móvel de 24 horas. O WAPE do melhor modelo é de 25\%, isto é, o erro médio é um quarto da concentração média observada.

A Tabela~\ref{tab:pareado_holdout} mostra as diferenças de MAE em relação à LSTM direta com perda de Huber e em relação à média móvel de 24 horas, as duas referências fixadas antes da execução.

\inputtabela{pareado_holdout}

Observe que o XGBoost erra menos que a LSTM direta em 0,212 (IC de \(-0,255\) a \(-0,169\)) e que as duas versões do Seq2Seq com atenção erram cerca de 0,10 menos; a LSTM recursiva e o Seq2Seq básico empatam com a LSTM direta. Todos os modelos treinados erram menos que a média móvel de 24 horas, com diferenças de \(-0,09\) (LSTM direta) a \(-0,30\) (XGBoost), com IC que exclui o zero.

As comparações entre todos os pares (Apêndice~\ref{ap:teste_historico}, Tabela~\ref{tab:matriz_holdout}) mostram um único empate com o melhor modelo: o XGBoost e a regressão Ridge têm diferença de \(-0,027\) (IC de \(-0,074\) a \(0,007\)), isto é, não é possível afirmar que um erra menos que o outro. Já o XGBoost erra menos que as duas versões do Seq2Seq com atenção (diferenças de 0,103 e 0,109), e a Ridge também erra menos que ambas (0,076 e 0,082, com IC que exclui o zero). As duas versões do Seq2Seq com atenção empatam entre si (diferença de \(-0,006\), IC de \(-0,028\) a \(0,013\)).

O \textit{ensemble} (média das previsões das quatro sementes; Tabela~\ref{tab:ensemble_holdout}) reduz o MAE em 0,01 a 0,06, conforme o modelo, mas não altera a ordenação: o XGBoost continua na frente (2,753). No teste histórico, o quadro é semelhante: a Ridge (2,673), o XGBoost (2,696) e o Seq2Seq com atenção adaptado (2,721) têm MAE próximos e sem diferença estatística entre si, e o Seq2Seq com atenção padrão fica atrás da Ridge.

\inputtabela{ensemble_holdout}

Em resumo, não há vencedor entre o XGBoost e a regressão linear, e nenhuma rede neural erra menos que eles. A escolha do modelo de aprendizado alterou o MAE em pouco mais de 7\%, uma diferença estatisticamente detectável, mas pequena.

\section{Referências Simples e Regressão Linear}
\label{sec:referencias_resultados}

As linhas dos previsores ingênuos e da Ridge na Tabela~\ref{tab:metricas_holdout} respondem se os modelos de aprendizado agregam valor. Sobre os previsores ingênuos, a resposta é sim: todos os modelos treinados têm MAE menor que o da melhor referência ingênua (a média móvel de 24 horas). A resposta é menos favorável quando a comparação é com a regressão Ridge, uma relação linear entre os 1.992 atributos e o valor a prever. Ela tem erro só 1\% maior que o do XGBoost, sem diferença estatística, e erra menos que todas as redes neurais treinadas. Portanto, nesta série, a relação entre o histórico e o valor futuro é bem capturada por um modelo linear, e os modelos não lineares (árvores e redes) não acrescentam ganho mensurável além dele.

\section{Erro por Horizonte}
\label{sec:erro_horizonte}

A Figura~\ref{fig:erro_por_horizonte_holdout} mostra o MAE em função do horizonte de previsão (1 a 24 horas). Observe que as linhas contínuas são os modelos treinados e as tracejadas, as referências.

\begin{figure}[!htbp]
    \centering
    \figuragerada{0.9\linewidth}{erro_por_horizonte_holdout}
    \caption{MAE por horizonte no período de confirmação (média das sementes). Linhas tracejadas: previsores de referência e regressão Ridge.}
    \label{fig:erro_por_horizonte_holdout}
\end{figure}

O erro cresce de forma lenta com o horizonte. Para o XGBoost, o MAE vai de 2,29 em 1 hora a 2,89 em 24 horas, um aumento de 26\%; para o Seq2Seq com atenção adaptado, de 2,61 a 2,98. No primeiro horizonte, o XGBoost erra 14\% menos que a persistência (2,29 contra 2,67), e a Ridge tem 2,40. A LSTM direta, por sua vez, tem erro quase constante (2,86 em 1 hora e 3,06 em 24 horas) e, nas primeiras horas, erra mais que a persistência: ao resumir a janela em um vetor único, ela parece perder precisão sobre o valor mais recente, o que os modelos que recebem o último valor observado no decodificador preservam. A LSTM recursiva e o Seq2Seq básico erram menos no início e mais no fim do horizonte (3,15 e 3,21 em 24 horas, contra 2,89 do XGBoost), o que é coerente com o acúmulo de erro típico das estratégias que realimentam a previsão. A ordenação entre os modelos no horizonte de 24 horas é a mesma da comparação global.

\section{Picos, Cauda e Suavização}
\label{sec:picos}

Modelos treinados com erro absoluto ou perda de Huber tendem a prever a mediana condicional \cite{gneiting2011making} e, por isso, produzem séries mais suaves que a observada e subestimam valores altos. A Tabela~\ref{tab:cauda_holdout} apresenta o MAE nas horas em que o PM2,5 observado excede 25, 35 e 50~\(\mu\)g/m\(^3\), a fração dessas horas em que o modelo subestima, a razão entre os desvios-padrão previsto e observado, e o percentil 99 da previsão (o observado é 33).

\inputtabela{cauda_holdout}

Todos os modelos subestimam de forma sistemática as horas de concentração alta. O MAE acima de 35~\(\mu\)g/m\(^3\) fica entre 17 e 21, isto é, cerca de metade do valor típico de um evento, e o modelo prevê valor abaixo do observado em 93\% a 100\% dessas horas. A razão entre os desvios-padrão fica entre 0,75 e 0,83 para os modelos treinados, e o percentil 99 previsto (27 a 30) fica abaixo do observado (33). A persistência e os previsores sazonais, que repetem valores passados, preservam a amplitude (razão próxima de 1), mas erram o momento dos picos.

A Tabela~\ref{tab:eventos_holdout} e a Figura~\ref{fig:janela_evento_holdout} descrevem os eventos, isto é, sequências contínuas de horas com PM2,5 acima do limiar. O período de confirmação tem 58 eventos acima de 35~\(\mu\)g/m\(^3\) e 243 acima de 25~\(\mu\)g/m\(^3\).

\inputtabela{eventos_holdout}

\begin{figure}[!htbp]
    \centering
    \figuragerada{0.78\linewidth}{janela_evento_holdout}
    \caption{Previsões (média das sementes) em torno do maior evento do período de confirmação, em painéis separados por modelo e com escala comum. Observe que a linha tracejada horizontal marca 35~\(\mu\)g/m\(^3\) e que as lacunas da série observada aparecem como quebras na linha.}
    \label{fig:janela_evento_holdout}
\end{figure}

No nível dos eventos, o MAE por evento acima de 35 fica entre 16,5 (Seq2Seq básico) e 19,0 (LSTM direta). A persistência (16,6) e a Ridge (16,8) estão entre os melhores, o que mostra que nenhum modelo treinado antecipa os picos melhor do que simplesmente repetir o último valor. Na comparação por evento com a LSTM direta (Tabela~\ref{tab:eventos_pareado_holdout}), todos os modelos treinados erram menos, mas a persistência também erra. As diferenças entre modelos nos picos são, portanto, pequenas diante do tamanho do erro de todos eles.

\inputtabela{eventos_pareado_holdout}

\subsection{Dinâmica Horária}
\label{sec:dinamica_horaria}

Preservar a amplitude não garante acertar o momento das variações. A Tabela~\ref{tab:dinamica_holdout} mostra a correlação entre os níveis previsto e observado, a correlação entre as variações de uma hora para a seguinte, e a fração de vezes em que a direção da mudança é acertada.

\inputtabela{dinamica_holdout}

A correlação entre os níveis é alta (0,78 a 0,81 nos modelos treinados), pois o modelo acompanha o nível médio. Porém, a correlação entre as variações horárias é baixa: de 0,13 a 0,17 nos modelos treinados, contra \(-0,33\) da persistência, e a direção da mudança é acertada em 54\% a 56\% das horas. O desvio-padrão das variações previstas (0,74 a 0,98) é muito menor que o das observadas (3,70). Assim, nenhum modelo acompanha as oscilações de hora em hora: eles acertam o nível diário, mas não o momento exato das variações.

\section{Efeito da Função de Perda}
\label{sec:efeito_perda}

Em etapas exploratórias anteriores, a vantagem do Seq2Seq adaptado em picos foi atribuída à combinação de arquitetura, perda ponderada e perda auxiliar. Como a perda ponderada foi aplicada a todas as arquiteturas, a Tabela~\ref{tab:fatorial_perda_holdout} separa o efeito da perda: para cada arquitetura, a diferença de MAE entre a perda ponderada e a de Huber, com intervalo de confiança pareado.

\inputtabela{fatorial_perda_holdout}

A perda ponderada não altera o MAE de quatro das cinco arquiteturas (empate estatístico). Em uma delas, o Seq2Seq adaptado, reduz o MAE em 0,049 (IC de \(-0,064\) a \(-0,032\)), o que equivale a 1,7\%. O efeito nos picos é pequeno: o MAE acima de 35 cai de 21,3 para 20,4 na LSTM direta, de 19,7 para 18,4 na recursiva, de 18,5 para 17,7 no Seq2Seq padrão e de 19,3 para 18,9 no adaptado (Tabela~\ref{tab:cauda_fatores_holdout}), e a razão entre os desvios-padrão sobe de 0,01 a 0,04. Portanto, a perda ponderada melhora modestamente os picos sem piorar o erro médio no período de confirmação. No teste histórico, ela piora o MAE da LSTM recursiva e do Seq2Seq básico, e portanto o efeito não é consistente entre arquiteturas e períodos.

\inputtabela{cauda_fatores_holdout}

\section{Fatores Adicionais e Desenho do Seq2Seq}
\label{sec:fatores_adicionais}

A Tabela~\ref{tab:extras_holdout} mostra o efeito da perda auxiliar de evento e do aprendizado residual, e a Tabela~\ref{tab:ablacao_desenho_holdout} remove, uma de cada vez, as quatro alterações que distinguem o Seq2Seq adaptado do padrão (Seção~\ref{sec:modelos}), mantendo os hiperparâmetros do modelo completo.

\inputtabela{extras_holdout}

\inputtabela{ablacao_desenho_holdout}

O aprendizado residual não ajuda: na LSTM direta não há diferença (+0,004, empate) e, no Seq2Seq adaptado, o MAE aumenta em 0,063 (IC de 0,018 a 0,112). Acrescentá-lo à combinação de perda ponderada e perda de evento piora em 0,099. Uma explicação possível é que a base sobre a qual o modelo aprende a correção, uma média de valores passados, é menos precisa que o próprio modelo, de modo que a correção precisa compensar o erro dela. A perda auxiliar de evento reduz o MAE em 0,021 sobre a perda ponderada (IC de \(-0,033\) a \(-0,011\)), um ganho pequeno.

Na ablação do desenho, remover o reinício do estado do decodificador, acrescentar o \textit{input feeding} ou remover o condicionamento da consulta ao horizonte não altera o MAE (empates). Remover as faixas de defasagem o aumenta em 0,015 (IC de 0,001 a 0,028), um efeito quase no limite da significância. No teste histórico, remover o reinício e o condicionamento da consulta chega a reduzir o MAE, com sinal oposto. Portanto, entre as quatro alterações, só as faixas de defasagem têm efeito distinguível de zero no período de confirmação, e ele é pequeno. Isso é coerente com o empate entre as duas versões do Seq2Seq com atenção observado na Seção~\ref{sec:comparacao_principal}.

\section{Variáveis Derivadas do PM10}
\label{sec:ablacao_pm10}

A Tabela~\ref{tab:ablacao_pm10_holdout} compara o conjunto completo de variáveis com duas remoções: o PM10 apenas no histórico, sem as duas variáveis derivadas dele para o futuro, e a remoção completa do PM10. Cada conjunto tem sua própria busca de hiperparâmetros, de modo que o conjunto completo não é favorecido.

\inputtabela{ablacao_pm10_holdout}

Para a LSTM direta, nenhuma remoção altera o MAE. Para o XGBoost, a remoção das variáveis derivadas aumenta o MAE em 0,010 (IC de 0,006 a 0,014), e a remoção do PM10 em 0,016 (IC de 0,010 a 0,022). Esses aumentos são estatisticamente distinguíveis de zero, mas equivalem a menos de 1\% do MAE e não têm importância prática. O PM10 medido pelo mesmo sensor traz, portanto, quase nenhuma informação adicional sobre o PM2,5 futuro.

\section{Interpretação da Atenção}
\label{sec:interpretacao_atencao}

Para testar se a atenção do Seq2Seq adaptado é funcionalmente relevante, aplicou-se a contraprova da Seção~\ref{sec:diagnostico_atencao}, sem retreinar. A Tabela~\ref{tab:atencao_seq2seq_attention_new__huber} mostra o erro quando os pesos aprendidos são substituídos por padrões fixos, e a Tabela~\ref{tab:atencao_diag_seq2seq_attention_new__huber} compara as medidas dos pesos aprendidos com as dos pesos uniformes.

\inputtabela{atencao_diagnostico_seq2seq_attention_new__huber}

\inputtabela{atencao_contrafactual_seq2seq_attention_new__huber}

Observe que zerar o contexto de atenção aumenta o MAE em 2,25 no período de confirmação (de 2,88 para 5,13): o modelo depende das informações que a atenção entrega. Já substituir os pesos por pesos uniformes sobre as 48 posições aumenta o erro em 0,15, e usar apenas a última hora, o mesmo horário do dia anterior ou o de dois dias antes o aumenta em 0,37, 0,48 e 0,76, respectivamente. O resultado mais informativo, porém, é a substituição por pesos uniformes \emph{dentro das seis faixas de defasagem}: o MAE aumenta em apenas 0,027 (1\%). As medidas dos pesos confirmam: a entropia normalizada dos pesos aprendidos (0,834) é quase igual à dos pesos uniformes por faixa (0,841), e as massas na última hora (0,207) e nas últimas seis horas (0,414 contra 0,417) também. Ou seja, a estrutura de seis faixas fornece quase todo o ganho, e a ponderação aprendida acrescenta pouco. Isso é coerente com o resultado da ablação, em que as faixas de defasagem são a única peça com efeito mensurável. Os pesos, porém, não devem ser lidos como uma explicação causal do modelo \cite{jain2019attention}: o que a contraprova mostra é o efeito no erro, não a razão pela qual o modelo prevê como prevê.

\section{Sensibilidades}
\label{sec:sensibilidades_resultados}

\subsection{Janela de Entrada}
\label{sec:sensibilidade_janela}

A Tabela~\ref{tab:janela_validacao} mostra o MAE de validação para janelas de 24, 48, 72 e 96 horas (o teste não participa desta escolha). Para a LSTM direta, a janela de 48 horas tem o menor erro (2,808), contra 2,842 com 24 horas, 2,832 com 72 e 2,829 com 96; as diferenças (0,02 a 0,03) são da ordem do desvio entre sementes (0,01 a 0,03). Para o XGBoost, o erro diminui com janelas maiores: 2,845 com 24 horas, 2,804 com 48, 2,787 com 72 e 2,772 com 96. A janela de 48 horas não é a ideal para o XGBoost, mas ganhar 0,03 (1\%) com uma janela de 96 horas é uma diferença pequena diante das demais deste capítulo, e não altera a ordenação entre as duas famílias na validação. A conclusão principal, portanto, não depende da janela na faixa testada.

\inputtabela{janela_validacao}

\subsection{Alvos Preenchidos, Busca de Hiperparâmetros e Vencedores}
\label{sec:sensibilidade_demais}

A Tabela~\ref{tab:sensibilidade_imputados_holdout} repete o cálculo do MAE incluindo, além das medições, os pares cujo alvo foi preenchido por interpolação. O erro sobre os alvos preenchidos é maior que o erro nos medidos para a maioria dos modelos (por exemplo, 3,54 contra 2,98 na LSTM direta); o XGBoost é exceção (2,51 contra 2,77). Isso mostra que o resultado sobre valores preenchidos depende de quanto o modelo se assemelha à interpolação, e não do fenômeno. A ordenação dos três melhores modelos se mantém (XGBoost, Seq2Seq padrão e adaptado); a LSTM recursiva e o Seq2Seq básico trocam de posição entre si. No teste histórico, a ordenação é idêntica. A conclusão, portanto, é robusta à exclusão dos alvos preenchidos.

\inputtabela{sensibilidade_imputados_holdout}

A Tabela~\ref{tab:verificacao_hpo_holdout} mostra o efeito de refazer a busca de hiperparâmetros para a perda ponderada, em vez de reaproveitar os do caso de Huber. Em ambas as arquiteturas, a diferença é estatisticamente indistinguível de zero (\(-0,004\) na LSTM direta e \(+0,006\) no Seq2Seq adaptado), o que indica que o reaproveitamento não distorce as conclusões sobre o efeito da perda.

\inputtabela{verificacao_hpo_holdout}

A Tabela~\ref{tab:vitorias_holdout} resume o melhor modelo em cada critério. O XGBoost é o melhor em todas as medidas globais e nos horizontes 1 e 24, empatado com a Ridge no MAE. O Seq2Seq básico é o melhor nas métricas de pico e no viés, com margens que não são estatisticamente relevantes diante do erro de todos os modelos nos picos.

\inputtabela{vitorias_holdout}

\section{O Que Limita a Qualidade da Previsão}
\label{sec:limites}

Três resultados apontam na mesma direção. Primeiro, modelos de famílias muito diferentes (regressão linear, árvores e redes com e sem atenção) terminam com MAE dentro de uma faixa estreita, de 2,77 a 2,98 no período de confirmação, e o modelo mais simples entre os treinados, a Ridge, empata com o melhor. Segundo, o erro já é alto na previsão de uma hora à frente: 2,29~\(\mu\)g/m\(^3\) para o melhor modelo, cerca de 21\% da concentração média (11,05), e o aumento até 24 horas é de apenas 26\%. Terceiro, o PM10 medido pela mesma estação pouco contribui (menos de 1\% do MAE), e nenhum modelo acompanha as oscilações de hora em hora ou os picos, que são subestimados em quase todas as horas.

Uma leitura compatível com esses resultados é que a informação contida no histórico de PM2,5 e PM10 e no calendário é o limite da previsão: a estação não mede variáveis meteorológicas, e estudos em outras cidades indicam que elas, como a altura da camada limite atmosférica, são as variáveis mais importantes para explicar o PM2,5 \cite{zanobini2024curitiba}. Esta é uma hipótese consistente com os dados, e não uma conclusão testada neste trabalho, pois testá-la exigiria variáveis meteorológicas que a estação não fornece. O que os resultados mostram, de forma direta, é que trocar de arquitetura ou de função de perda altera o erro em poucos por cento, e que nenhuma delas resolve a imprevisibilidade dos picos e das oscilações horárias.

\section{Ameaças à Validade}
\label{sec:ameacas}

Há limitações que condicionam a leitura dos resultados. Primeiro, o estudo usa uma única estação, o que limita a generalização espacial. Segundo, o teste histórico foi consultado em etapas anteriores e não é confirmatório; o período de confirmação é a única evidência independente, e cobre apenas três anos, com nível médio de PM2,5 variável entre eles. Terceiro, os fatores que reaproveitam os hiperparâmetros do caso de Huber podem subestimar o potencial de cada variante, embora a verificação da Seção~\ref{sec:sensibilidade_demais} não tenha encontrado viés em duas arquiteturas. Quarto, quatro sementes fornecem pouca precisão para inferência entre famílias; o bootstrap incorpora a incerteza da amostra de teste, mas não a de novas inicializações além das quatro usadas. Quinto, o número de eventos de pico é limitado (58 acima de 35~\(\mu\)g/m\(^3\) no período de confirmação), e quase não há valores acima de 35 em 2023. Sexto, a validação tem cerca de 32\% de ausências e uma lacuna de 48 dias preenchida por repetição do último valor, o que pode enfraquecer a parada antecipada e a escolha do ponto de controle. Sétimo, a janela de 48 horas não é a ideal para o XGBoost, como mostrado na Seção~\ref{sec:sensibilidade_janela}. Por fim, o limiar de 35~\(\mu\)g/m\(^3\), os pesos da perda ponderada e as faixas de defasagem são decisões de projeto e não resultados estabelecidos na literatura.
